\documentclass[12pt]{article}
%% BUGS:
%% - The \begin{problem} ... \end{problem} is only used in exams??

\usepackage{url, graphicx, epstopdf}

% page layout
\usepackage[letterpaper, textheight=9.5in, textwidth=5.5in]{geometry}
\usepackage{setspace}
\setstretch{1.08}
\pagestyle{plain}
\sloppy\sloppypar\raggedbottom\frenchspacing\thispagestyle{empty}

% problem set header
\newcommand{\psname}{Problem Set}
\newcommand{\startps}[1]{\section*{Computational Physics --- \psname~{#1}}}

% problem formatting
\usepackage{amsthm}
\theoremstyle{definition}
\newcommand{\problemname}{Problem}
\newtheorem{problem}{\problemname}
\newcommand{\probref}[1]{\problemname~\ref{#1}}
\theoremstyle{remark}
\newcommand{\notename}{Note}
\newtheorem{note}{\notename}[problem]
\newcommand{\noteref}[1]{\textit{\notename}~\ref{#1}}
\newcounter{subpart}[problem]
\newcommand{\subpart}{\par\addvspace{1ex}\refstepcounter{subpart}\textsl{(\alph{subpart})}~}

% words
\newcommand{\foreign}[1]{\textsl{#1}}
\newcommand{\vs}{\foreign{vs}}

% code
\newcommand{\code}[1]{\texttt{\detokenize{#1}}}
\usepackage{xcolor}
% Define custom colors for code highlighting
\definecolor{codegreen}{rgb}{0,0.6,0}
\definecolor{codegray}{rgb}{0.5,0.5,0.5}
\definecolor{codepurple}{rgb}{0.58,0,0.82}
\definecolor{backcolour}{rgb}{0.99,0.99,0.99}
\usepackage{listings}
\lstdefinestyle{pythonstyle}{
    backgroundcolor=\color{backcolour},   
    commentstyle=\color{codegreen},
    keywordstyle=\color{magenta},
    numberstyle=\tiny\color{codegray},
    stringstyle=\color{codepurple},
    tabsize=4
}
\lstset{
    basicstyle=\small\ttfamily,
    columns=fixed,
    breaklines=false,
    keepspaces=true,
    showspaces=false,
    showstringspaces=false,
    showtabs=false,
    showlines=*,
    style=pythonstyle,
    language=Python,
    literate={~}{{\url{~}}}1
}

% math
\usepackage{amsmath}
\newcommand{\dd}{\mathrm{d}}
\newcommand{\e}{\mathrm{e}}

% primary units
\newcommand{\rad}{\mathrm{rad}}
\newcommand{\kg}{\mathrm{kg}}
\newcommand{\m}{\mathrm{m}}
\newcommand{\s}{\mathrm{s}}

% secondary units
\renewcommand{\deg}{\mathrm{deg}}
\newcommand{\km}{\mathrm{km}}
\newcommand{\cm}{\mathrm{cm}}
\newcommand{\mm}{\mathrm{mm}}
\newcommand{\ft}{\mathrm{ft}}
\newcommand{\mi}{\mathrm{mi}}
\newcommand{\AU}{\mathrm{AU}}
\newcommand{\ns}{\mathrm{ns}}
\newcommand{\h}{\mathrm{h}}
\newcommand{\yr}{\mathrm{yr}}
\newcommand{\N}{\mathrm{N}}
\newcommand{\J}{\mathrm{J}}
\newcommand{\eV}{\mathrm{eV}}
\newcommand{\W}{\mathrm{W}}
\newcommand{\Pa}{\mathrm{Pa}}

% derived units
\newcommand{\mps}{\m\,\s^{-1}}
\newcommand{\mph}{\mi\,\h^{-1}}
\newcommand{\mpss}{\m\,\s^{-2}}

% random units
\newcommand{\au}{\mathrm{a.u.}}


\begin{document}

\startps{2}
Due Friday 2026 September 18 at 18:00 on Brightspace.
Be sure to explicitly give credit where it is due, on all problems.
Also credit your buddies and LLMs (as per course policy).
Failure to give explicit credits will result in points deducted.

\begin{problem}
  \subpart Read in the data file at \href{https://drive.google.com/uc?export=download&id=1YCUGRTV5uSZeWPHa7CNXzn-YB6OC7pnm}{this link}; it contains one very long column of data.
  Now sum the entries to get one total sum of all the numbers in the file.
  Now re-order the entries and sum again. Did you get the same number?
  Now consider very specific, deterministic orderings that might be better or worse, and sum those.

  \subpart If you get different answers in different orderings, which one is the correct sum?
  How can you verify that one sum is correct and the others are incorrect?

  \subpart Do your sums manually (with a dumb \code{for} loop maybe?) and also do them with \code{numpy.sum()}.
  What can you conclude about how numpy does its sums? Anything?
\end{problem}

\begin{note}
  You will have to figure out how to read CSV files or equivalent text files.
  There is a standard Python implementation, and there is also an \code{astropy} implementation.
\end{note}

\begin{note}
  Hogg abhors text files. Why? What's wrong with them?
\end{note}

\begin{problem}
  Re-do \psname~1, \problemname~1(a), but with a root-finding algorithm.
  Write a simple implementation of bisection root finding, but which operates on a function $f(x)$ that returns \code{True} if a condition on $x$ is true and \code{False} if that condition is false.
  If you don't know what the bisection method is, look it up on \textit{Wikipedia}.
  The root-finding code should take in three arguments: The function $f$, a value $x_1$ that is on the ``true'' side of the condition, and a value $x_2$ that is on the ``false'' side of the condition.
  Here's pseudo-code for your binary function:
\begin{lstlisting}
def f(x):
    # Note that x is the natural log of delta
    return 1. + np.exp(x) > 1.
\end{lstlisting}
  The calling sequence for your root-finding function should look something like
\begin{lstlisting}
def find_root(f, x1, x2, tol=1e-6):
    if f(x1):
        xtrue, xfalse = x1, x2
    else:
        xtrue, xfalse = x2, x1
    assert f(xtrue) & ~f(xfalse)
    # [write your iterative bisection algorithm here]
    return x
\end{lstlisting}
Fill in that function and run it and give the output.
Is it consistent with what you got for \psname~1?
Write your code ``verbose'' such that it shows the internal values for \code{xtrue} and \code{xfalse} at each iteration.
\end{problem}

\begin{note}
  I wrote my code snippet to be agnostic about which of $x_1,x_2$ is smaller or larger, and which is true and which is false.
  All the code requires is that one of them be true and one of them be false.
  See if you can make your code agnostic in this way, and then test that it is.
\end{note}

\begin{note}
  I wrote my code snippet such that if it is called badly, it \emph{fails}.
  This is good practice for scientific programming.
\end{note}

\begin{problem}
Newman, \textit{Computational Physics, 2ed}, Problem 4.3.
\end{problem}

\begin{note}
  If you don't have the second edition of the textbook---or if you don't have the textbook at all---then consult the PDF file of Chapter~4 that was shared on the Discord.
\end{note}

\begin{problem}
  Newman, \textit{Computational Physics, 2ed}, Problem 4.5, except make the output of the problem a plot:
  Plot the (absolute value of the) inaccuracy of your answer vs the number of samples $N$, and also vs wall-clock time $T$.
  Extend your plot to a few seconds of wall-clock time, or as far as you can stand it.
  Make your plots log--log, and make sure you obey all of Hogg's plotting guidance.
\end{problem}

\begin{note}
  You might have to learn how to use the Python \code{time} package.
\end{note}

\begin{note}
  Hogg's plotting guidance was given in Lecture on 2026-09-08.
  If you missed that Lecture, consult your buddies!
\end{note}

\end{document}
